\hypertarget{huawei-ascend-computing-platform-6.1}{%
\subsubsection{1. Huawei Ascend Computing Platform
(6.1)}\label{huawei-ascend-computing-platform-6.1}}

\textbf{What I learned:}

\begin{itemize}
\tightlist
\item
  The Ascend processor uses the \textbf{Da Vinci architecture} , which
  is specifically designed for AI computing with \textbf{99.2\% of its
  core dedicated to Cube units} for matrix multiplication. A single Cube
  unit can perform 4096 FP16 operations per clock cycle (16x16 matrix
  multiply).
\item
  The AI Core contains three execution units: \textbf{Cube} (matrix),
  \textbf{Vector} (vector), and \textbf{Scalar} (scalar/program
  control). The storage system includes input/output buffers and
  registers to reduce bus traffic.
\item
  AI processors are classified into \textbf{training} (high throughput,
  large memory, e.g., Atlas 900) and \textbf{inference} (low power, low
  latency, e.g., Atlas 200/300 series).
\item
  The \textbf{Atlas computing platform} scales from edge (Atlas 200,
  500) to data center (Atlas 800 training/inference, Atlas 900 cluster).
  For instance, Atlas 900 can deliver \textbf{20.4 PFLOPS FP16} with
  ultra‑high energy efficiency.
\item
  Huawei uses Ascend in real‑world projects, like the \textbf{Peng Cheng
  Cloud Brain II} , reducing TCO by 9.3\% with the same computing power,
  and the \textbf{Traffic Brain} , improving violation detection by 10×.
\end{itemize}

\textbf{Why it was interesting/useful:}

\begin{itemize}
\tightlist
\item
  I was impressed by the \textbf{Cube unit} 's massive parallelism (up
  to 4096 operations per cycle) and how tightly the hardware is
  optimized for tensor operations. The comparison with GPUs (only a
  small fraction of Tensor Cores) clearly showed why dedicated NPUs
  achieve higher efficiency.
\item
  The \textbf{hardware‑software co‑design} (CANN, DVPP, runtime manager)
  demonstrates how a full stack enables high performance. The real‑world
  case of reducing power and space by 60--80\% in data centers makes the
  technology tangible.
\end{itemize}

\begin{center}\rule{0.5\linewidth}{0.5pt}\end{center}

\hypertarget{huawei-cloud-ei-platform-6.2}{%
\subsubsection{2. Huawei Cloud EI Platform
(6.2)}\label{huawei-cloud-ei-platform-6.2}}

\textbf{What I learned:}

\begin{itemize}
\tightlist
\item
  \textbf{ModelArts} is a one‑stop AI development platform that covers
  the entire ML lifecycle: data labeling, training (ExeML for
  beginners), deployment, and MLOps. The \textbf{ExeML engine} allows
  people without programming skills to build models in three steps.
\item
  Advanced services include \textbf{OptVerse AI Solver} for operations
  optimization, \textbf{federated learning} to break data silos, and
  \textbf{ModelBox} for device‑edge‑cloud joint inference, improving E2E
  performance by up to 3×.
\item
  General AI capabilities: \textbf{vision services} (100+ table types,
  face search \textgreater95.5\%), \textbf{NLP} (industry‑leading
  accuracy in NER, sentiment analysis, summarization), \textbf{knowledge
  graph} (full‑lifecycle management), and \textbf{Conversational Bot
  Service} (handled workload of 179 human agents annually at Huawei).
\item
  The platform also offers third‑party integration and an \textbf{AI
  Gallery} bridging supply and demand.
\end{itemize}

\textbf{Why it was interesting/useful:}

\begin{itemize}
\tightlist
\item
  I really liked \textbf{ExeML} because it shows how AI can be
  democratised---users only need to upload and label data, train, then
  publish a model. The \textbf{federated learning} capability is crucial
  for industries like healthcare and finance that cannot share raw data.
\item
  The \textbf{CBS bot} case study (serving 350,000+ Huawei staff, 65\%
  manual substitution rate) proved the immediate ROI of cloud AI
  services. I can see myself using such cloud APIs for rapid
  prototyping.
\end{itemize}

\begin{center}\rule{0.5\linewidth}{0.5pt}\end{center}

\hypertarget{huawei-device-ai-platforms-6.3}{%
\subsubsection{3. Huawei Device AI Platforms
(6.3)}\label{huawei-device-ai-platforms-6.3}}

\textbf{What I learned:}

\begin{itemize}
\tightlist
\item
  \textbf{HarmonyOS} is a distributed OS enabling seamless collaboration
  across devices.
\item
  \textbf{HMS Core} offers capabilities in seven domains to build an
  intelligent app ecosystem.
\item
  \textbf{ML Kit} provides a wide range of AI APIs (text, language,
  image, face/body) that work across Android, iOS, and HarmonyOS, with
  on‑device processing for data security.
\item
  \textbf{HiAI} is a three‑layer open platform: \textbf{Foundation}
  opens NPU acceleration (150+ operators), \textbf{Engine} delivers rich
  AI capabilities out of the box, and \textbf{Service} connects users to
  intelligent services.
\item
  \textbf{MindSpore Lite} is an ultra‑fast, lightweight AI engine that
  supports CPU/GPU/NPU, model compression, and can run on everything
  from phones to IoT devices. It is compatible with MindSpore,
  TensorFlow Lite, Caffe, and ONNX.
\end{itemize}

\textbf{Why it was interesting/useful:}

\begin{itemize}
\tightlist
\item
  The concept of \textbf{on‑device AI} (ML Kit, HiAI) is fascinating
  because it ensures privacy (data never leaves the device) and reduces
  latency. As a developer, I can use the \textbf{HiAI Foundation} 's
  operators to get NPU acceleration without deep hardware knowledge.
\item
  \textbf{MindSpore Lite} 's claim of ``all‑scenario support'' and fast
  deployment for extreme environments makes it very practical for edge
  AI projects. The fact that it can quantize and compress models for
  tiny devices is something I want to explore further.
\end{itemize}

\begin{center}\rule{0.5\linewidth}{0.5pt}\end{center}

\hypertarget{cuttingedge-ai-applications-7}{%
\subsubsection{4. Cutting‑edge AI Applications
(7)}\label{cuttingedge-ai-applications-7}}

\textbf{What I learned:}

\begin{itemize}
\tightlist
\item
  \textbf{Reinforcement learning (RL):} agents learn optimal policies
  through interaction with the environment. Algorithms can be direct
  (policy optimization) or indirect (value‑based).
\item
  \textbf{GANs:} a generator and a discriminator compete, eventually
  creating realistic fake data. Applications include image generation,
  text‑to‑image, super‑resolution, and speech enhancement.
\item
  \textbf{Knowledge graphs:} structured semantic networks of entities,
  properties, and relationships. Construction involves extraction,
  fusion, and storage. Applications: search engines, Q\&A bots,
  recommendation systems, and domain‑specific solutions (e.g., oil \&
  gas).
\item
  \textbf{Intelligent driving:} history from Carnegie Mellon's Navlab
  (1995, 98.1\% autonomous across the US) to DARPA challenges
  (2004--2007) that fostered modern self‑driving cars. Today's
  autonomous driving cloud service offers a one‑stop platform with data,
  training, and simulation services.
\end{itemize}

\textbf{Why it was interesting/useful:}

\begin{itemize}
\tightlist
\item
  The \textbf{GAN explanation} made the adversarial concept very clear:
  ``the generator creates fakes, the discriminator tries to spot them.''
  I found the practical uses (resolution enhancement, text‑to‑image)
  inspiring.
\item
  The \textbf{historical timeline of intelligent driving} was
  eye‑opening---it highlighted that many foundational ideas date back to
  the 1990s. The DARPA challenges were a brilliant catalyst for the
  industry. The modular autonomous driving cloud platform
  (data‑training‑simulation) gave me insight into how one might build
  such a system today.
\end{itemize}

\begin{center}\rule{0.5\linewidth}{0.5pt}\end{center}

\hypertarget{quantum-computing-and-machine-learning-8}{%
\subsubsection{5. Quantum Computing and Machine Learning
(8)}\label{quantum-computing-and-machine-learning-8}}

\textbf{What I learned:}

\begin{itemize}
\tightlist
\item
  \textbf{Quantum mechanics basics:} superposition, entanglement, and
  uncertainty. A \textbf{qubit} can be in a superposition of \textbar0⟩
  and \textbar1⟩, enabling exponential state space.
\item
  \textbf{Key algorithms:}

  \begin{itemize}
  \tightlist
  \item
    \emph{Deutsch algorithm} solves a simple problem in one query
    instead of two, demonstrating quantum advantage.
  \item
    \emph{Shor} factorizes integers in polynomial time, threatening RSA.
  \item
    \emph{Grover} searches unsorted databases with $\sqrt{N}$ complexity.
  \item
    \emph{HHL} solves linear systems exponentially faster than classical
    methods.
  \end{itemize}
\item
  \textbf{Quantum machine learning (QML)} combines classical
  data/algorithms with quantum resources. Variational algorithms (e.g.,
  VQE for chemistry) use a hybrid classical‑quantum loop to find ground
  state energies. QML promises acceleration for tasks like PCA, SVM, and
  reinforcement learning.
\item
  \textbf{MindSpore Quantum} is an open‑source framework that integrates
  with MindSpore, offering high‑performance simulators (up to 30+
  qubits), rich tutorials, and seamless cloud deployment. I got hands‑on
  with a code snippet that builds a circuit, computes expectation
  values, and gradients.
\end{itemize}

\textbf{Why it was interesting/useful:}

\begin{itemize}
\tightlist
\item
  The idea of \textbf{exponential parallelism} through superposition is
  mind‑blowing---the Deutsch algorithm showing we can solve a problem
  with one function call instead of two illustrated a tangible, though
  small, quantum advantage.
\item
  \textbf{MindSpore Quantum} 's simplicity (PIP install, visualization,
  built‑in libraries) makes it approachable even for an ML student like
  me. The example of calculating expected values and gradients in just a
  few lines of code makes me excited to try quantum machine learning
  experiments. The VQE application to chemistry shows how quantum
  computing can address real scientific challenges today.
\end{itemize}
